\section{Preliminaries: Energy Matching and JKO }

Energy matching learns a scalar potential \( V_\theta:\mathcal{X}\to\mathbb{R} \), parameterized by \(\theta\), on a continuous data space \(\mathcal{X}\). It simultaneously serves two complementary roles: (i)~off-manifold guidance, directing samples toward the data manifold via the gradient field \(-\nabla_x V_\theta\), and (ii)~near-equilibrium refinement, characterizing the data distribution on the manifold as an \gls{ebm} with Gibbs density \(\propto \exp(-V_\theta(x))\).

These roles separate model learning into a \emph{global} transport-alignment and a \emph{local} density-refinement term. A simple time-dependent temperature schedule to switch between these regimes, e.g., \(\epsilon(t)=0\) for \(t<1\) (transport) and \(\epsilon(t)=1\) for \(t\ge 1\) (mixing), (i) injects exploration, (ii) mitigates mode collapse, and (iii) stabilizes sampling from the Gibbs distribution \(\propto \exp(-V_\theta(x))\).

Energy Matching can be understood through an optimization-in-probability-spaces perspective, expressed via the \gls{jko}) scheme from density $\rho_t$ to $\rho_{t+\Delta t}$:
\begin{align}
\label{jko_em}
\rho_{t+\Delta t}
&=
\arg\min_{\rho}\;
\frac{1}{2\Delta t}\,\inf_{\gamma\in\Gamma(\rho_t,\rho)}\mathbb{E}_{(x,y)\sim\gamma}[\,c(x,y)\,]\nonumber\\
&\quad+\;\int V_\theta(y)\,\rho(y)\,\mathrm{d}y\nonumber\\
&\quad+\;\epsilon(t)\int \rho(y)\log\rho(y)\,\mathrm{d}y,
\end{align}
where the cost function \( c(x,y) \) measures the amount of displacement when transporting mass from \(x\) to \(y\), and constitutes a design choice---in continuous spaces typically chosen as \( c(x,y)=\|x - y\|^2 \).

The coupling \(\gamma\in\Gamma(\rho,\rho_t)\) is a joint distribution that characterizes a transport plan: it precisely describes how mass from the current configuration \(\rho_t\) is relocated to form the new distribution \(\rho\). Intuitively, \(\gamma(x,y)\) encodes ``how much mass from each location \(x\sim\rho_t\) is transported to each new location \(y\sim\rho\)'', thus determining the optimal reassignment of mass that minimizes the expected transportation cost \(\mathbb{E}_{(x,y)\sim\gamma}[\,c(x,y)\,]\).

\paragraph{Off-manifold transport guidance.}
Based on first-order optimality conditions in Wasserstein space~\citep{lanzetti2024variational,lanzetti2025firstorder} applied to the functional in~\cref{jko_em}, Energy Matching defines two clear learning objectives. Given tractably coupled pairs \((x,y)\sim\gamma\) (e.g., obtained via minibatch matching rather than solving for the globally optimal coupling), we associate \(x_0\) with the noise/source distribution (\(\rho_{0}\)) and \(x_1\) with the data/target distribution (\(\rho_{\mathrm{data}}\)). Intermediate points are sampled as \(x_t=(1-t)x_0+t x_1\) for \(t\sim\Unif([0,1])\). Defining the empirical transport direction \(v=x_1 - x_0\), the \(V_\theta\) is trained via a flow-like objective:
\begin{equation}
\mathcal{L}_{\mathrm{Flow}}
=
\mathbb{E}_{(x,y)\sim\gamma,\,t\sim\Unif([0,1])}\big[\|\nabla_{x_t} V_\theta(x_t) + v\|^2\big].
\end{equation}

\paragraph{Near-equilibrium energy refinement.}
At equilibrium (\(t\to\infty\) in Eq.~\eqref{jko_em} and in the mixing regime \(\epsilon(t)=1\)), the stationary distribution induced by \(V_\theta\) is given by:
\begin{equation}
\rho_{t\to\infty}(x)\propto \exp(-V_\theta(x)).
\end{equation}

Refinement of the equilibrium density is achieved using the contrastive loss (CL) (see Appendix \ref{app:energy-refinement}):
\begin{equation}
\label{eq:cl}
\mathcal{L}_{\mathrm{CL}}
=
\mathbb{E}_{x^{+}\sim \rho_{\mathrm{data}}}[V_\theta(x^{+})]
-
\mathbb{E}_{x^{-}\sim\operatorname{sg}(\rho_{t\to\infty})}[V_\theta(x^{-})],
\end{equation}
where negative samples \(x^{-}\) are generated via Langevin dynamics initialized either from the source or from data, with gradients stopped (\(\operatorname{sg}\)) during sampling. Training begins with an initial warm-up phase using only \(\mathcal{L}_{\mathrm{Flow}}\), followed by joint optimization of:
\begin{equation}
\mathcal{L}(\theta)=\mathcal{L}_{\mathrm{Flow}}(\theta)+\lambda_{\mathrm{CL}}\,\mathcal{L}_{\mathrm{CL}}(\theta).
\end{equation}

\paragraph{Sampling (two regimes).}
Sampling starts from a transport-aligned drift and switches to refinement via an effective temperature
$\epsilon(t)$. In continuous time, this corresponds to the time-inhomogeneous \gls{sde}
\begin{equation}
\label{em_sde}
\mathrm{d}x_t
=
-\nabla_{x_t} V_\theta(x_t)\,\mathrm{d}t
+\sqrt{2\,\epsilon(t)}\,\mathrm{d}W_t,
\end{equation}
which reduces to an \gls{ode} when $\epsilon(t)=0$ and becomes stochastic as $\epsilon(t)\to 1$.
In practice, sampling corresponds to explicitly integrating \cref{em_sde} over a fixed time interval.

\subsection{Regressor-guided Diffusion}
Both DeFoG~\citep{qinmadeira2024defog} and Variational Flow Matching (VFM)~\citep{eijkelboom2024vfm,eijkelboom2025controlledvfm} achieve controllable graph diffusion by conditioning transitions on noise-dependent regressors predicting desired properties. Formally, these regressors approximate conditional distributions of graph properties $y$ given intermediate noisy graph states $x_t$, i.e., $\rho(y\mid x_t,t)$ or $\rho_t(y\mid x_t)$. Further details and empirical analyses illustrating how prediction accuracy behaves at different noise levels are provided in Appendix~\ref{app:regressor-noise}. {\color{red}\textbf{TODO: move to appendix/baselines.}}
